% 第 7 章：配电网潮流（前推回代 / 调压器外环 / PV 功率曲线 / 支路回算）
% 本章文件由 \input 引入，不含导言区。
\section{配电网潮流}\label{sec:dist}\label{sec:distribution}

\subsection{功能定位与入口}

\compmeta{analysis::solve\_distribution\_pf}{\srcpath{include/hacdcpf/power\_flow/distribution\_power\_flow.hpp}}

配电网潮流族以辐射状馈线的\textbf{前推回代}（Backward-Forward Sweep, BFS；
Shirmohammadi et al.\ 1988）为主路径，面向单相（正序）等值模型
（include/hacdcpf/power\_flow/distribution\_power\_flow.hpp）。两个重载入口（:solve\_distribution\_pf）：
\begin{quote}\footnotesize
\texttt{DPFResult solve\_distribution\_pf(const ACSystem\& ac\_sys, const DPFOptions\& opt = \{\});}\\
\texttt{DPFResult solve\_distribution\_pf(const HybridPowerSystem\& sys, const DPFOptions\& opt = \{\});}
\end{quote}
实现位于 \srcpath{src/power\_flow/distribution\_power\_flow.cpp}（单相内核
:390--540；入口 :1041--1387）。结果索引空间与 \fld{ac\_sys.buses}/\fld{branches}
\textbf{完全同序}（hpp:75--76, 432--433）：\fld{vm\_pu[i]} 对应
\texttt{buses[i]}，\fld{p\_branch\_mw[b]} 对应 \texttt{branches[b]}。
\texttt{HybridPowerSystem} 重载先经 \srcpath{projection::RichToCanonicalOperator::apply}
投影到 canonical 模型（cpp:1067--1068），求解后把 \fld{vm\_pu}/\fld{va\_deg}
按 \fld{bus\_merge\_map} 以强度量语义广播回原始母线编号空间
（:1048--1057，回投影契约见第~\ref{sec:overview}~章；注意
\fld{p\_branch\_mw} 等支路量\textbf{不}回投影，仍按 canonical 支路序索引）。

如实说明定位边界：\srcpath{solve\_distribution\_pf} 是库级 API，
\textbf{未}接入公共门面 \srcpath{include/hacdcpf/api/hacdcpf.hpp}，也未接入 GUI
后端 \srcpath{tests/run\_gui\_server.cpp} 的 \texttt{method} 分发；环网与多电源
三相场景应改用三相 NR（第~\ref{sec:threephase}~章）或统一 Newton
（第~\ref{sec:newton}~章）。同文件还提供 abc 域三相 BFS
\srcpath{solve\_three\_phase\_distribution\_pf}（hpp:456--464，
实现 cpp:1632--1804），其公式与 NR 对照细节见第~\ref{sec:threephase}~章，
本章仅在支路回算一处引用其损耗口径。

\subsection{辐射拓扑处理与 BFS 树}
\srcpath{build\_bfs\_tree}（cpp:74--117）从首个 SLACK 母线出发对在役支路
邻接表做广度优先遍历，生成三个并行数组：BFS 访问序 \fld{order}、父母线
\fld{parent\_bus}、父支路 \fld{parent\_branch} 及父侧方向标志
\fld{parent\_is\_from\_side}（:build\_bfs\_tree）。无 SLACK 直接抛异常（:solve\_distribution\_pf\_ac\_snapshot）。
诚实口径：
\begin{itemize}
  \item \textbf{辐射假设不强制校验}：头文件仅要求调用方先用
        \srcpath{topology\_analysis.hpp} 的 \texttt{is\_radial()} 自查
        （hpp:429--430）；环网送入时 BFS 树按先访问者确定父支路，环路支路
        不参与扫描但仍出现在支路潮流回算中，结果\textbf{不报错也不保证正确}；
  \item \textbf{与根不连通的孤岛}：未进入 \fld{order} 的母线电压保持初值
        （\fld{vm\_pu}、相角 $0$），不参与残差计算，结果中静默保留（:solve\_distribution\_pf\_ac\_snapshot）；
  \item \textbf{多 SLACK}：只取首个 SLACK 为根（:solve\_distribution\_pf\_ac\_snapshot），其余 SLACK 母线
        当作普通节点处理，其出力由 \fld{compute\_net\_injection} 按普通发电机
        扣减（:260--267 仅跳过根）。
\end{itemize}

\subsection{前推回代算法}
净注入 \fld{compute\_net\_injection}（cpp:199--312）按
$S^{\mathrm{net}}_i=(S^{\mathrm{load}}_i-S^{\mathrm{gen}}_i)/S_{\mathrm{base}}$
汇总：优先取 \fld{loads} 表（为空时回退母线级 \fld{pd\_mw}/\fld{qd\_mvar}，
:208--225）、充电站（kW/kvar 换算，:228--235）、可选母线与可分组并联补偿
（\fld{include\_shunts}，:238--257）、常规/静态/可再生发电机、PV（经功率曲线，
:290--299，见本节下文）与储能（放电为正，:301--309），发电机侧一律跳过根节点。
基准容量 \fld{base\_mva}$\le 0$ 时回退 $100$~MVA（:solve\_distribution\_pf\_ac\_snapshot）。注意 ZIP 负荷在此
被\textbf{恒功率近似}（注释 :214--216），不随电压迭代更新。

内层扫描 \srcpath{run\_bfs\_sweep}（:run\_bfs\_sweep）。记复变比
$t=|\mathrm{tap}|\angle\delta$（$|\mathrm{tap}|<10^{-12}$ 视为 $1$，:121--125），
支路串联电流
\[
I_s=\big(V_f/t - V_t\big)/z,\qquad z=r+jx,\quad |z|<10^{-20}\Rightarrow I_s=0
\quad(:127\text{--}135).
\]
\textbf{回推}（:run\_bfs\_sweep）：由 BFS 序逐点计算注入电流
$I_b=\overline{S^{\mathrm{net}}_j/V_j}$（$|V_j|\le10^{-12}$ 时取 $0$），再逆序
把子支路电流累加进父支路，跨越变比时按父侧方向共轭折算：父在 from 侧
$I_{\mathrm{up}}=I_{\mathrm{child}}/\overline{t}$，父在 to 侧
$I_{\mathrm{up}}=\overline{t}\,I_{\mathrm{child}}$（:branch\_subtree\_current\_seen\_from\_parent）。
\textbf{前推}（:branch\_child\_voltage\_from\_parent）沿 BFS 序更新电压：
\[
V_j=\begin{cases}
V_i/t - z\,I_b & \text{父在 from 侧}\\[2pt]
t\,\big(V_i - z\,\overline{t}\,I_b\big) & \text{父在 to 侧}
\end{cases}
\qquad(:158\text{--}168).
\]
收敛判据为全母线 $\max|\Delta V|<\texttt{opt.tol}$（:run\_bfs\_sweep）；初值取各母线
\fld{vm\_pu}、相角 $0$ 的平启动（:solve\_distribution\_pf\_ac\_snapshot）。收敛后由
\srcpath{compute\_single\_phase\_branch\_flow}（:compute\_single\_phase\_branch\_flow）回算两端功率
$S_f=V_f\overline{I_f}\,S_{\mathrm{base}}$、$S_t$ 与损耗
$S_{\mathrm{loss}}=S_f+S_t$，写入 \fld{p\_branch\_mw}/\fld{q\_branch\_mvar}
与总损耗（:solve\_distribution\_pf\_ac\_snapshot）。

可选\textbf{无功限值外环}（\fld{enforce\_q\_limits}，:438--507）：内层收敛后
逐台发电机复算母线所需无功 $Q_{\mathrm{req}}$（由其相连支路潮流求和再扣回本地
负荷，:457--476），越 \fld{qmax\_mvar}/\fld{qmin\_mvar} 则把 \fld{qg} 钉到限值、
标记已钳位并重启平启动；外环上限 \fld{q\_limit\_max\_outer\_iter}=30 次，
\fld{qmax}=\fld{qmin}=0 视为无限值跳过（:solve\_distribution\_pf\_ac\_snapshot）。钳位不可逆、无迟滞恢复，
与统一 Newton 的 PV$\leftrightarrow$PQ 外环（第~\ref{sec:newton}~章）相比更简化。

\begin{paramtable}
\fld{max\_iter} & $N_{\mathrm{it}}$ & 次 & $50\sim500$ & 100 & BFS 内层迭代上限（hpp:65）\\
\fld{max\_control\_iter} & $N_{\mathrm{ctl}}$ & 次 & $10\sim100$ & 100 & 调压器离散控制外环上限（:DPFOptions）\\
\fld{tol} & $\varepsilon$ & pu & $10^{-10}\sim10^{-6}$ & $10^{-6}$ & 收敛容差，$\max|\Delta V|$（:DPFOptions）\\
\fld{verbose} & --- & bool & --- & false & 打印逐迭代残差（:DPFOptions）\\
\fld{include\_shunts} & --- & bool & --- & true & 计入母线/补偿器 $g_s,b_s$（:DPFOptions）\\
\fld{enforce\_q\_limits} & --- & bool & --- & false & 启用发电机 Q 限值外环（:DPFOptions）\\
\fld{q\_limit\_max\_outer\_iter} & --- & 次 & $10\sim30$ & 30 & Q 限值外环上限（:DPFOptions）\\
\end{paramtable}

\subsection{调压器（OLTC/稳压器）离散控制外环}
\texttt{HybridPowerSystem} 重载检测到任一 \fld{regulator\_controls} 启用即进入
控制外环（cpp:1058--1066, 1106--1380）：每轮先投影、跑一趟 BFS，再逐个调压器
评估一次档位动作，全部 \texttt{within\_band} 才判 \fld{control\_converged}。
控制律为 OpenDSS RegControl 风格（伏特域，非标幺）：

\begin{enumerate}
  \item \textbf{监控电压}：默认监控 \fld{winding} 侧端母线，\fld{monitored\_bus}$\neq0$
        时取远端母线（:solve\_distribution\_pf）；电压幅值乘 $U_{\mathrm{base}}=$
        \fld{base\_kv}$\times1000$（:629--635, 1196--1202，缺 \fld{base\_kv} 抛异常）；
        PT 变比远端时优先 \fld{remote\_ptratio}（:solve\_distribution\_pf）；
  \item \textbf{线路压降补偿（LDC）}：\fld{r\_volts}/\fld{x\_volts} 非零时
        要求 \fld{ct\_primary\_amps}$>0$，取调压绕组侧支路电流
        $I_w$（安培，基值 $I_{\mathrm{base}}=S_{\mathrm{base}}\!\times\!10^{3}/U_{\mathrm{base,kV}}$，
        :1238--1241），补偿相量
        $\Delta U_c=(I_w/I_{\mathrm{ct}})\,(R+jX)$（:solve\_distribution\_pf）；
  \item \textbf{控制电压与死区}：
        $U_{\mathrm{ctl}}=\big|V_{\mathrm{mon}}/\mathrm{ptratio}+\Delta U_c\big|$
        （:solve\_distribution\_pf）；$U_{\mathrm{ctl}}<U_{\mathrm{reg}}-B/2-10^{-6}$ 升压、
        $>U_{\mathrm{reg}}+B/2+10^{-6}$ 降压，否则 \texttt{within\_band} 停
        （:1248--1250, 1270--1299；$10^{-6}$~V 数值迟滞）；
  \item \textbf{档位动作}：连续变比
        $\mathrm{tap}_{\mathrm{pu}}=1+(\mathrm{tap\_pos}-\mathrm{tap\_neutral})\cdot s\%/100$
        （:transformer\_tap\_pu）；单步动作
        $\Delta\mathrm{pos}=\pm\min(\texttt{max\_tap\_change},|\mathrm{tap\_max}-\mathrm{tap\_min}|)$，
        方向由 \fld{tap\_side} 决定（升压时 $=1$ 取 $+1$、$=0$ 取 $-1$，:567--569），
        并钳位到 $[\mathrm{tap\_min},\mathrm{tap\_max}]$（:solve\_distribution\_pf）；触限即
        fail-close（\texttt{tap\_limit\_reached\_*}）；
  \item \textbf{振荡防护与诚实停机}：以全体启用调压器的 \fld{tap\_pos} 向量
        做历史状态集，重复即 \texttt{cycle\_detected}；耗尽
        \fld{max\_control\_iter} 记 \texttt{max\_control\_iter\_reached}；两种情形
        均置 \fld{control\_converged}=\fld{converged}=false（:solve\_distribution\_pf）。
\end{enumerate}
不支持的模式\textbf{显式 fail-close} 而非静默近似（:1164--1187, 1309--1317）：
\fld{reversible}、\fld{monitored\_node}$\neq1$、\fld{tap\_winding} 与
\fld{tap\_side} 不匹配、远端监控与 LDC 同用、缺失离散档位网格等均抛出原因串
写入 \fld{stop\_reason}。逐轮决策轨迹见 \fld{regulator\_trace}
（hpp:78--96），终态汇总 \fld{regulator\_states}（hpp:98--128）。

\begin{paramtable}
\fld{vreg\_volts} & $U_{\mathrm{reg}}$ & V & $115\sim125$ & 0 & 目标控制电压（include/hacdcpf/model/ac\_components.hpp:RegulatorControl）\\
\fld{band\_volts} & $B$ & V & $1\sim4$ & 0 & 全死区宽度，半宽 $B/2$（:RegulatorControl）\\
\fld{ptratio} & $n_{\mathrm{PT}}$ & --- & $20\sim120$ & 0 & PT 变比；$\le10^{-12}$ 抛错（:RegulatorControl）\\
\fld{remote\_ptratio} & --- & --- & 同左 & 0 & 远端 PT 变比；$\le0$ 回退 \fld{ptratio}（:RegulatorControl）\\
\fld{ct\_primary\_amps} & $I_{\mathrm{ct}}$ & A & $100\sim1000$ & 0 & CT 一次电流，LDC 必填（:RegulatorControl）\\
\fld{r\_volts}/\fld{x\_volts} & $R+jX$ & V & $0\sim30$ & 0 & LDC 补偿阻抗（伏特域，:1127--1128）\\
\fld{max\_tap\_change} & --- & 档 & $1\sim4$ & 1 & 单轮最大动作档数；$<1$ 抛错（:RegulatorControl）\\
\fld{winding}/\fld{tap\_winding} & --- & 1|2 & --- & 0 & 监控/调档绕组（:RegulatorControl）\\
\fld{monitored\_bus} & --- & id & --- & 0 & 0=监控绕组端母线（:RegulatorControl）\\
\fld{reversible} & --- & bool & --- & false & 逆功率模式；当前不支持（:RegulatorControl）\\
\end{paramtable}

\subsection{PV 功率曲线}
\compmeta{powerflow::compute\_pv\_power\_mw}{\srcpath{include/hacdcpf/power\_flow/pv\_power\_curve.hpp}}

经验单二极管族曲线（纯头文件实现，include/hacdcpf/power\_flow/pv\_power\_curve.hpp:pv\_module\_current）：
\[
I(V)=I_{\mathrm{sc}}\Big(1-\big(V/V_{\mathrm{oc}}\big)^{a}\Big)^{b}
      \Big(1-c\big(V/V_{\mathrm{mpp}}\big)^{2}\Big),
\qquad a=10.0,\ b=0.547596,\ c=0.023812 .
\]
阵列出力（:compute\_pv\_power\_mw）在 MPPT 点取值（$V=V_{\mathrm{mpp}}$）：
\[
P_{\mathrm{MW}}=I(V_{\mathrm{mpp}})\,V_{\mathrm{mpp}}\,
  N_{\mathrm{series}}N_{\mathrm{parallel}}\,
  \frac{G}{1000}\,\eta_{\mathrm{inv}}\times10^{-6}.
\]
\fld{voc}/\fld{isc}/\fld{vmpp} 任一 $\le0$ 时回退静态 \fld{p\_mw}（:compute\_pv\_power\_mw）。
相关字段默认值（\srcpath{include/hacdcpf/model/ac\_components.hpp:PVSystem}）：
\fld{inverter\_eff}=0.97、\fld{irradiance}=1000~W/m$^2$、串并联数
\fld{num\_series}/\fld{num\_parallel} 默认 0（触发回退）。BFS 在净注入阶段调用
（src/power\_flow/distribution\_power\_flow.cpp:compute\_net\_injection），GUI 时序仿真亦复用同一函数
（\srcpath{tests/run\_gui\_server.cpp:has\_profile（lambda）}）。诚实口径：该模型是
\textbf{MPPT 单点估值}而非 $P(V)$ 电压依赖注入——潮流迭代中 PV 出力恒定，
曲线自变量并不取端电压；\fld{alpha\_isc}/\fld{beta\_voc} 温度系数字段存在
但此函数未使用。

\subsection{支路功率回算与结果结构}
两套支路回算并存，勿混淆：
\begin{itemize}
  \item \textbf{BFS 内部回算} \srcpath{compute\_single\_phase\_branch\_flow}
        （src/power\_flow/distribution\_power\_flow.cpp:compute\_single\_phase\_branch\_flow）：仅串联阻抗 + 复变比，
        $S_f=V_f\overline{I_f}S_{\mathrm{base}}$，损耗 $=S_f+S_t$；
        不含对地支路充电 $b/2$——BFS 结果中的 \fld{q\_branch\_mvar} 与统一
        Newton 口径在高充电线路上存在系统性差异；
  \item \textbf{统一 Newton 回算} \srcpath{powerflow::compute\_branch\_flows}
        （\srcpath{src/power\_flow/branch\_flow.cpp:compute\_branch\_flows}，规范头
        \srcpath{include/hacdcpf/power\_flow/assembly/branch\_flow.hpp}，
        \srcpath{power\_flow/branch\_flow.hpp} 为转发头）：MATPOWER $\pi$ 模型
        $y_{ff}=(y_s+jb/2)/|t|^2$、$y_{ft}=-y_s/\overline{t}$、
        $y_{tf}=-y_s/t$（:compute\_branch\_flows），$|z|=0$ 与停运支路跳过，
        供 \srcpath{solve\_power\_flow} 派生结果使用（见第~\ref{sec:newton}~章）。
\end{itemize}
输出 \fld{DPFResult}（include/hacdcpf/power\_flow/distribution\_power\_flow.hpp:DPFResult）：
\fld{vm\_pu}/\fld{va\_deg}（相角单位\textbf{度}）、\fld{p\_branch\_mw}/
\fld{q\_branch\_mvar}（送端 from 侧）、总损耗、\fld{converged}/
\fld{control\_converged}、\fld{iterations}、\fld{residual}（末轮
$\max|\Delta V|$），以及调压器 \fld{regulator\_states}/\fld{regulator\_trace}。
\fld{specialized\_path\_id} 记录三绕组三角等专属致密 NR 试验路径的触发与
fail-close 原因（cpp:637--648, 995--996）。三相 BFS 的损耗回算保留
$Z_{abc}$ 互耦交叉项（注释明确反对仅用自阻抗 $|I_p|^2$ 近似，
src/power\_flow/distribution\_power\_flow.cpp:solve\_three\_phase\_distribution\_pf）。

\subsection{验证与校核}
全部 BFS/调压器对照集中在 \srcpath{tests/test\_opendss\_pf\_compare.cpp}
（Catch2 + OpenDSS 桥，受 \srcpath{HACDCPF\_ENABLE\_OPENDSS\_COMPARE} 门控，
\srcpath{tests/CMakeLists.txt:254--257}，OpenDSS 缺失时整组 skip）：
\begin{itemize}
  \item \textbf{3 母线 1ph 基线}（:TEST\_CASE）：vs OpenDSS 快照，
        $|V|$ 容差 $10^{-3}$、相角 $0.1^\circ$、送端 $P/Q$ 各 $10^{-3}$、
        总损耗 $10^{-3}$；并联补偿变体（:TEST\_CASE）与物理电容器 vs BFS 恒功率
        并联近似的探索性对照（:TEST\_CASE）；
  \item \textbf{case33bw 全馈线}（:TEST\_CASE）：Baran \& Wu 33 节点平衡三相
        OpenDSS 模型 vs BFS 单相等值，33 母线 $|V|$ 容差 $2\times10^{-3}$、
        相角 $0.2^\circ$，32 支路三相合计 $P/Q$ 各 $10^{-3}$，总损耗
        $5\times10^{-3}$；同测试含 OpenDSS 自相一致性回归（相间幅值差
        $<10^{-4}$、$\pm120^\circ$ 偏差 $<10^{-3}$~deg）；
  \item \textbf{调压器闭环}（:TEST\_CASE）：本地 PT、本地 PT+LDC、
        RemotePTRatio 三算例，终态档位数与 OpenDSS \texttt{regcontrol}
        精确一致（容差 $10^{-12}$），控制电压与按 DSS 数据独立复算值差
        $<0.2$~V（:verify\_regulator\_case\_matches\_repository），终止原因 \texttt{within\_band}；另有运行时
        语义矩阵（:1970，remote/local/LDC 组合逐场景核对 \fld{stop\_reason}
        与 DSS 桥一致）与带母线合并的失败路径索引回归（:2490, :2519）；
  \item \textbf{固定档变压器}：1ph 直通变压器（:TEST\_CASE）、固定档 vs 变比感知
        BFS（:TEST\_CASE）、LV 侧固定档端到端（:TEST\_CASE）。
\end{itemize}
诚实边界：单相 BFS \textbf{没有}独立单元测试文件，全部覆盖经由上述
OpenDSS 门控集成测试；环网/多 SLACK/死岛行为无测试；case33bw 对照为
\textbf{平衡}系统（单相等值严格成立的特例）；IEEE 13 节点正式交叉校核
当前状态为 \texttt{deferred}、原因 \texttt{model\_limitations}（:3235--3264,
:3438），即不平衡全馈线的 BFS 精度尚未验收——不平衡场景以三相 NR 的
回归（见第~\ref{sec:threephase}~章）为准。
